\section{Results}\label{sec:7}

We report results in the order of the principles of \secref{sec:3}, mapping each measured property to a user-facing interpretation.
All numbers come from committed benchmark artifacts; the supplementary material reports the underlying tables, the calibration curves, and the full method-comparison matrix.

\leadpara{Recommendation quality (G1 and G4)}
\tabref{tab:progression} reports \ndcg{}, precision@5, and recall@5 for the recommendation baselines, the portal-default composite, and the strongest single configuration.
The portal-default composite---the strategy the user actually sees in the prototype---reaches \ndcg{} $= 0.949$, with P@5 $= 0.303$ and R@5 $= 0.983$.
The strongest single configuration, a multimodal soft-embedding ranker with skill weights set to $0.7$, reaches \ndcg{} $= 0.969$ with P@5 $= 0.313$ and R@5 $= 1.000$.
A learned logistic-regression fusion ties the best progression run at \ndcg{} $= 0.968$; fixed multimodal blending and hierarchical skill weighting both reach $0.924$.
Lexical baselines remain close (BM25 \ndcg{} $= 0.901$, TF--IDF $0.905$); pure semantic cosine is $0.911$; a richer semantic template nudges to $0.922$.
A cross-encoder reranker reaches \ndcg{} $= 0.939$ but lowers the score from the bi-encoder baseline by $0.108$ on this corpus, consistent with the prior report's decision to leave it disabled by default.

\begin{center}
\refstepcounter{table}\label{tab:progression}
\par\smallskip\noindent\textbf{Table~\thetable:} Progression of \ndcg{}, P@5, and R@5 across recommendation baselines. The portal-default composite is the strategy the user sees; the multimodal soft-embed configuration is the strongest single configuration. Differences between rows should be read with the corpus size caveat from \secref{sec:6} in mind.
\par\smallskip
\begin{tabular}{@{}lrrr@{}}
\toprule
Configuration & P@5 & R@5 & nDCG@5 \\
\midrule
Multimodal soft embed (skill weight $0.7$)  & 0.313 & 1.000 & 0.969 \\
Soft embed + cross-encoder rerank            & 0.307 & 0.983 & 0.939 \\
RRF ensemble of four list views              & 0.293 & 0.950 & 0.935 \\
Multimodal Jaccard (skill weight $0.7$)      & 0.293 & 0.950 & 0.933 \\
Semantic cosine, rich template               & 0.300 & 0.950 & 0.922 \\
Semantic cosine                              & 0.280 & 0.900 & 0.911 \\
TF--IDF (lexical)                            & 0.307 & 0.983 & 0.905 \\
BM25 (lexical)                               & 0.307 & 0.983 & 0.901 \\
\midrule
Portal-default composite                     & 0.303 & 0.983 & 0.949 \\
\bottomrule
\end{tabular}
\par\vspace{0.4em}
\footnotesize{Source: committed benchmark artifact \texttt{paper\_progression\_summary.json}. Gated by a regression test with a $\pm 0.04$ \ndcg{} tolerance.}
\end{center}

For the user, the practical interpretation is that the top-five list on average places the most relevant jobs within reach: recall@5 saturates at $1.000$ for the best configuration, meaning every resume retrieves all of its labeled relevant jobs within the top five.
The precision@5 of approximately $0.31$ means that, on this corpus, roughly three of every ten top-five slots are labeled relevant; the user should still expect to see plausible but not relevant jobs in the remainder and should rely on the explanation drawer and the calibrated confidence to identify them.

\leadpara{Statistical significance}
Paired bootstrap tests ($N{=}5000$, seed $42$) confirm the hybrid gains are not noise.
The portal-default composite significantly improves \ndcg{} over semantic-only ranking ($\Delta = +0.071$, $p = 0.048$).
Multimodal blending and reciprocal rank fusion both significantly beat semantic cosine ($p = 0.010$ and $p = 0.009$ respectively).
The lexical baselines do not differ significantly from semantic cosine ($p > 0.25$), consistent with their close scores on this small corpus.
For the user, the significance result is a guardrail: a user who notices the explanation panel weighting both text and skills is not seeing a quirk of the corpus.

\leadpara{Explanation faithfulness (G2)}
\tabref{tab:explainability} reports the explainability metrics for the rule-based explainer (the deployed default) and the LLM-template explainer (an offline comparison) on the 150 composite top-5 matches.
The rule-based explainer in the prototype achieves faithfulness $0.745$, specificity $0.627$, and consistency $1.000$ across the synthetic profile variants.
Skill-mention coverage is $25.3\%$, indicating that the rule-based explainer names a concrete matched or missing skill in roughly a quarter of explanation bullets; the remainder are general statements that the explainer is honest about not having a specific reason to assert.
For the user, the practical interpretation is that a user who reads an explanation can trust the named components to be the components that produced the score, and can read the absence of a specific skill as a signal of low specificity rather than as a hidden model behavior.
The supplementary material reports the same metrics for the LLM-based explainer; the LLM-template explainer reaches similar faithfulness (0.747) and slightly higher skill-mention coverage (0.267), at the cost of a marginally higher hallucination rate (2.7\% vs 2.0\%); both explainers have 100\% component alignment with the channel scores that produced the rank.

\begin{center}
\refstepcounter{table}\label{tab:explainability}
\par\smallskip\noindent\textbf{Table~\thetable:} Explainability evaluation summary, rule-based explainer (deployed) vs.\ LLM-template explainer (offline comparison). Faithfulness is the pass rate on the automated checks (skill mention, no hallucination, component alignment). Consistency is the bullet-set Jaccard on synthetic similar-profile pairs. Source: committed benchmark artifact \texttt{explainability\_report.json}.
\par\smallskip
\footnotesize
\setlength{\tabcolsep}{3pt}
\begin{tabular*}{\linewidth}{@{\extracolsep{\fill}}lcccccc@{}}
\toprule
Explainer & Faithful. & Specif. & Skill (\%) & No halluc. (\%) & Comp. align (\%) & Consist. \\
\midrule
Rule-based (deployed)       & 0.745 & 0.627 & 25.3 & 98.0 & 100.0 & 1.000 \\
LLM-template (comparison)   & 0.747 & 0.633 & 26.7 & 97.3 & 100.0 & 1.000 \\
\bottomrule
\end{tabular*}
\par\vspace{0.4em}
\normalsize
\end{center}

\leadpara{Confidence calibration (G3)}
\tabref{tab:calibration} reports the calibration metrics before and after Platt scaling.
The uncalibrated composite has \ece{} $= 0.40$ on the held-out calibration set, which is too poorly calibrated to expose directly to the user.
After Platt scaling with parameters $a = 0.298$ and $b = -2.116$, the \ece{} drops to $0.032$, a reduction of approximately an order of magnitude.
The Brier score of the calibrated system is $0.093$.
For the user, the practical interpretation is that a confidence value displayed next to a ranked item is a usable signal: a confidence of $0.9$ is associated with a top-rank outcome approximately nine times out of ten, and a confidence of $0.5$ is associated with one approximately half the time.
The calibration is committed and reproducible; the parameters are exposed in the prototype.

\begin{center}
\refstepcounter{table}\label{tab:calibration}
\par\smallskip\noindent\textbf{Table~\thetable:} Calibration metrics before and after Platt scaling. The Platt parameters are learned by maximum likelihood on the held-out calibration set of 21 strong + 26 partial labels. ECE is the expected calibration error; Brier is the mean squared error of the confidence prediction. Source: committed benchmark artifact \texttt{calibration.json}.
\par\smallskip
\begin{tabular}{@{}lcccc@{}}
\toprule
Configuration & $a$ & $b$ & ECE & Brier \\
\midrule
Uncalibrated composite (Platt removed) & -- & -- & 0.400 & 0.245 \\
Calibrated composite (Platt, MLE) & 0.298 & $-2.116$ & 0.032 & 0.093 \\
\bottomrule
\end{tabular}
\par\vspace{0.4em}
\end{center}

\leadpara{Counterfactual probe (G3)}
\tabref{tab:counterfactual} reports the per-pair breakdown of the 10-pair counterfactual probe.
Seven of the ten counterfactual pairs in the probe are flagged as having a single-field edit that would change the rank by at least one position; in all seven cases, the top-ranked item remains stable, and the maximum score shift across all ten pairs is $0.017$.
The probe is small but supports the design choice to expose counterfactual previews in the user interface: when the system names a component as a reason, a single-field edit to that component is usually sufficient to move the rank.
For the user, the practical interpretation is that the counterfactual view of \figref{fig:portal}(d) is not a toy: the predicted effect of an edit corresponds to an actual effect on the next ranking.

\begin{center}
\refstepcounter{table}\label{tab:counterfactual}
\par\smallskip\noindent\textbf{Table~\thetable:} Counterfactual probe on 10 synthetic demographic-pair edits. The top-1 item is stable for all 10 pairs; the top-5 overlap is the number of items in common with the unedited profile's top-5. A pair is flagged when the score delta exceeds 0.005 or the rank delta exceeds 1. No protected attributes are inferred from real users. Source: committed benchmark artifact \texttt{fairness\_eval.json}.
\par\smallskip
\scriptsize
\begin{tabular*}{\linewidth}{@{\extracolsep{\fill}}lccccc@{}}
\toprule
Pair & Top-1 stable & Top-5 overlap & Max rank $\Delta$ & Max score $\Delta$ & Flagged \\
\midrule
name\_gender\_01     & Yes & 5/5 & 2 & 0.0067 & No  \\
name\_gender\_02     & Yes & 5/5 & 1 & 0.0050 & Yes \\
name\_ethnicity\_01  & Yes & 5/5 & 0 & 0.0110 & Yes \\
name\_ethnicity\_02  & Yes & 4/5 & 3 & 0.0167 & Yes \\
nationality\_01      & Yes & 5/5 & 1 & 0.0061 & Yes \\
nationality\_02      & Yes & 5/5 & 1 & 0.0077 & No  \\
hometown\_01         & Yes & 5/5 & 1 & 0.0082 & Yes \\
hometown\_02         & Yes & 4/5 & 2 & 0.0100 & Yes \\
pronouns\_01         & Yes & 5/5 & 0 & 0.0029 & No  \\
email\_domain\_01    & Yes & 5/5 & 1 & 0.0057 & Yes \\
\midrule
\textbf{Summary}   & \textbf{10/10} & \textbf{9.0/10} avg & \textbf{1.2} avg & \textbf{0.0080} avg & \textbf{7/10} flagged \\
\bottomrule
\end{tabular*}
\par\vspace{0.4em}
\normalsize
\end{center}

\leadpara{Adversarial probe (fairness and robustness)}
\tabref{tab:fairness} reports the selection-rate disparate-impact ratios on synthetic proxy groups.
The disparate-impact ratio on the experience tier dimension is $0.82$; on the remote-preference dimension it is $0.75$ (parity is $1.0$).
Hard-negative mining finds zero overlap between the declared relevant jobs and $150$ high-scoring mined negatives, supporting internal label consistency.
The probe is narrow and is reported as such; we discuss what it does and does not establish in \secref{sec:8}.

\begin{center}
\refstepcounter{table}\label{tab:fairness}
\par\smallskip\noindent\textbf{Table~\thetable:} Adversarial fairness probe: selection-rate disparate-impact ratios (DIR) on synthetic proxy groups. Parity is 1.0. The experience-tier probe groups candidates by self-reported seniority (mid vs senior); the remote-preference probe groups by declared remote-work preference (remote-seeker vs onsite-flexible). Source: committed benchmark artifact \texttt{fairness\_eval.json}.
\par\smallskip
\scriptsize
\setlength{\tabcolsep}{3pt}
\begin{tabular*}{\linewidth}{@{\extracolsep{\fill}}lcc@{}}
\toprule
Dimension & Selection rate (group 1) & Selection rate (group 2) \\
\midrule
Experience tier (mid vs senior)  & 0.522 & 0.429 \\
Remote preference (remote vs onsite-flexible) & 0.450 & 0.600 \\
\midrule
Dimension & Disparate-impact ratio (DIR) & Parity \\
\midrule
Experience tier  & 0.82 & 1.0 (parity) \\
Remote preference & 0.75 & 1.0 (parity) \\
\bottomrule
\end{tabular*}
\par\vspace{0.4em}
\normalsize
\end{center}

\leadpara{Latency}
\tabref{tab:latency} reports the nDCG@5 and per-query latency for the main scoring methods on the 15-job pool.
Lightweight scorers complete in well under one millisecond per resume on the $15$-job pool; the full composite bi-encoder run averages approximately $0.4$ ms per query.
The cross-encoder reranker averages approximately $141.7$ ms per query, which is approximately $340$ times slower than the bi-encoder on this setup and is the reason it remains disabled by default.
For the user, the practical interpretation is that the prototype responds to a snapshot confirmation in well under one perceptual latency unit, so the system's feedback feels immediate and the user can iterate on edits without a perceptible wait.

\begin{center}
\refstepcounter{table}\label{tab:latency}
\par\smallskip\noindent\textbf{Table~\thetable:} Quality--latency trade-off across retrieval methods on the 15-job pool with $K = 5$. The soft-skill embed reranker is approximately 100$\times$ slower than the composite bi-encoder and is reported as a research variant, not the deployed default. Source: committed benchmark artifact \texttt{phase11\_summary.csv}.
\par\smallskip
\scriptsize
\setlength{\tabcolsep}{3pt}
\begin{tabular*}{\linewidth}{@{\extracolsep{\fill}}lrrl@{}}
\toprule
Method & nDCG@5 & Latency (ms/q) & Notes \\
\midrule
TF--IDF (lexical)         & 0.905 & 0.05  & Lexical baseline \\
BM25 (lexical)            & 0.902 & 0.09  & Lexical baseline \\
Exact skill overlap       & 0.748 & 0.17  & Lexical baseline \\
Skills Jaccard             & 0.748 & 0.17  & Skill overlap \\
Semantic (L2)              & 0.878 & 0.21  & Embedding \\
Semantic (cosine)          & 0.878 & 0.23  & Embedding \\
Multimodal (cosine, $w=0.7$) & 0.924 & 0.51  & Embedding, deployed default \\
RRF ensemble (4 lists)     & 0.913 & 1.35  & Embedding, fusion \\
Soft-skill embed reranker  & 0.869 & 51.81 & Embedding, research variant \\
Cross-encoder reranker     & 0.939 & 141.7 & Embedding, disabled by default \\
\bottomrule
\end{tabular*}
\par\vspace{0.4em}
\normalsize
\end{center}

\leadpara{Engineering surface}
The prototype is implemented as a small web stack; the candidate, employer, and admin portals sit above a shared gateway.
The repository contains $302$ Python unit tests and $39$ Node tests, for a total of $341$ tests, with a continuous-integration gate that fails the build on a \ndcg{} regression greater than $\pm 0.04$ against the committed benchmark.
The full engineering detail, including endpoint maps, default request parameters, and the test surface, is reported in the supplementary material.
The engineering surface is reported because the user-facing properties depend on it: a prototype whose tests gate on a \ndcg{} tolerance is a prototype whose user-facing claims are reproducible.
